\chapter{Solenoidal image and resolution of Navier--Stokes}

\section{Hydrodynamic domain and Galerkin publication}

We fix
\[
 H^s_\sigma(\mathbb T^3)=
 \{u\in H^s(\mathbb T^3;\mathbb R^3):
   \nabla\!\cdot u=0,\ \int_{\mathbb T^3}u=0\},
 \qquad s>\frac52,\quad\nu>0,
\]
with
$u_0\in C^\infty\cap H^s_\sigma$ and
$f\in C^\infty([0,\infty);H^\infty_\sigma)
\cap L^2_{\rm loc}([0,\infty);H^{s-1}_\sigma)$. For the Leray--Galerkin projection $P_M$ set
\[
 A_M=-P_M\Delta,
 \qquad B_M(v,w)=P_MP_\sigma((v\cdot\nabla)w),
 \qquad N_M(u):=B_M(u,u),
\]
and the finite field satisfies
\begin{equation}\label{eq:pdf2-ns-galerkin}
 \partial_tu_M+\nu A_Mu_M+B_M(u_M,u_M)=f_M,
 \qquad u_M(0)=P_Mu_0.
\end{equation}
The equation \eqref{eq:pdf2-ns-galerkin} is a publication of the state, not its
complete definition.

So that the notation does not change halfway through the chain, the
four energy functionals of order $s$ are fixed once and for all:
\begin{equation}\label{eq:pdf2-ns-energy-functionals}
 \begin{aligned}
 \mathscr E_{s,M}(u)&:=\frac12\|\Lambda^su\|_{L^2}^2,
 &\mathscr D_{s,M}(u)&:=\|A_M^{1/2}\Lambda^su\|_{L^2}^2,\\
 \mathscr B_{s,M}(u)&:=
   \operatorname{Re}\langle\Lambda^sN_M(u),\Lambda^su\rangle,
 &\mathscr F_{s,M}(u,f)&:=
   \operatorname{Re}\langle\Lambda^sP_Mf,\Lambda^su\rangle.
 \end{aligned}
\end{equation}
The additional subscript $j$ denotes the same four functionals evaluated on
the published coordinate of depth level $j$; it introduces no other
functionals.
The Galerkin equation implies exactly
\begin{equation}\label{eq:pdf2-ns-energy-identity}
 \frac d{dt}\mathscr E_{s,M}(u_M)
 +\nu\mathscr D_{s,M}(u_M)+\mathscr B_{s,M}(u_M)
 =\mathscr F_{s,M}(u_M,f).
\end{equation}

\section{Source object and coordinates that are indeed constructed}

The solenoidal image is not added a posteriori. It is exactly the autonomous
construction $\mathfrak G_{\rm sol}$ produced before this chapter and published in
\eqref{eq:pdf2-g-sol-explicito}; it preserves the
Galerkin state, refinement maps, oriented sectors, dissipative operator,
projected advection, carry, counit $E_9$, inclusion $J_9$,
projections $P_9=J_9E_9$ and $Q_{\rm micro}=I-P_9$, memory, route and terminal.
In particular,
\[
 E_9J_9=I,
 \qquad P_9^2=P_9,
 \qquad Q_{\rm micro}=I-P_9.
\]
These identities type the decomposition between macroscopic publication and
microscopic memory; by themselves they do not constitute a global bound.

At nonadic depth $j$, the state retains
\begin{equation}\label{eq:pdf2-ns-estado}
 \Xi_{M,j}^{\rm NS}=
 \bigl(u_M;q_{A,M,j},q_{V,M,j},T_{M,j};
 \begin{gathered}
 \text{triads, sheet, orientation, carry, boundary,}\\
 \text{route, archive and memory}
 \end{gathered}
 \bigr).
\end{equation}
The enriched return combines all branches before evaluation:
\begin{equation}\label{eq:pdf2-ns-retorno}
 \mathcal R_{M,J,t}^{\rm enr}
 =\mathcal R_{1,M}\mathfrak U_{M,J,t}^*,
 \qquad
 \mathcal R_{M,J,t}^{\rm enr}Z_{M,J,t}=u_M(t).
\end{equation}
An isolated norm branch $3^J$ is not the return.

For $N>M$, the commutation defect with the nonlinearity is the carry
\begin{equation}\label{eq:pdf2-ns-carry}
 C_{M\leftarrow N}:=
 P_MB(u_N,u_N)-B_M(P_Mu_N,P_Mu_N).
\end{equation}
Setting it equal to zero without proving commutation removes precisely the high-frequency
interaction that the refinement must preserve.

The two positive memories remain separate. With
$r=\pi_{\rm HMT}/\varphi_{\rm HMT}^2>1$,
\begin{equation}\label{eq:pdf2-ns-dh}
 D_M=(r-1)(\mathcal M_M^+-\mathcal M_M^-),
 \qquad
 H_M=(1+r)(\mathcal M_M^++\mathcal M_M^-),
\end{equation}
and
\[
 \dot D_M=\mathscr B_{s,M},
 \qquad
 \dot H_M=\frac{1+r}{r-1}|\mathscr B_{s,M}|.
\]
$D_M$ publishes sign; $H_M$ preserves positive magnitude. Merging them destroys the
orientation or the positivity.

\section{Complete PC--Fock tower and terminal storage}

The elevation does not truncate at degree two:
\begin{equation}\label{eq:pdf2-ns-pcf}
 \mathcal V_M^{\rm PCF}(u)
 :=\bigoplus_{n\ge0}\frac{R_{\rm F}^n}{\sqrt{n!}}u^{\odot n},
 \qquad
 \|\mathcal V_M^{\rm PCF}(u)\|^2
 =e^{R_{\rm F}^2\|u\|_{H^s}^2}.
\end{equation}
The radius is part of the coordinate and cannot disappear from the reader. For
any $\rho>0$ one fixes
\begin{equation}\label{eq:pdf2-ns-normalized-degree-one-reader}
 \rho^{[\rho]}_1:=\rho^{-1}\pi_1:
 \Gamma_s^\rho(V)\longrightarrow V,
 \qquad
 \rho^{[\rho]}_1\mathcal V_\rho^{\rm PCF}(u)=u.
\end{equation}
The initial level uses $\rho^{[R_{\rm F}]}_1$ and level $j$ uses
$\rho^{[r_j]}_1$; this normalization is maintained when depth is withdrawn.
On the coherent diagonal, $z_{n,M}=u_M^{\odot n}$ satisfies
\begin{equation}\label{eq:pdf2-ns-carleman}
 \dot z_{n,M}=L_{n,M}z_{n,M}
 +\mathcal N_{n,n+1;M}z_{n+1,M}
 +\mathcal F_{n,n-1;M}(f_M)z_{n-1,M}.
\end{equation}
First-degree and second-degree readers are
\[
 z_M(u)=\nu^{1/2}A_M^{1/2}\Lambda^su,
 \qquad
 w_M(u)=\nu^{-1/2}A_M^{-1/2}\Lambda^sN_M(u),
\]
and Young gives
\begin{equation}\label{eq:pdf2-ns-young-port}
 |\mathscr B_{s,M}(u)|\le
 \beta\nu\mathscr D_{s,M}(u)+\frac{\|w_M(u)\|^2}{4\beta}.
\end{equation}
The uniform estimate of the Fourier rows controls this quadratic reader in
the initial state. It does not yet control its future propagations.

Let
\[
 \mathcal N_{M,j}:=
 \frac{r-1}{r+1}H_{M,j}-D_{M,j}^{\rm mem}
 =2(r-1)M_{M,j}^-\ge0.
\]
Over a horizon $T$, the terminal storage is
\begin{equation}\label{eq:pdf2-ns-storage}
 \begin{aligned}
 G_{M,j,T}(t)
 &=\mathscr E_{s,M,j}(t)+2(r-1)(M^-_{M,j}(T)-M^-_{M,j}(t)),\\
 \dot G_{M,j,T}+\nu\mathscr D_{s,M,j}+|\mathscr B_{s,M,j}|
 &=\mathscr F_{s,M,j}.
 \end{aligned}
\end{equation}
To type the terminal, the sheet coordinate is preserved in the complete solenoidal state.
\[
 q^-_{M,j}:=\sqrt{2(r-1)\dot M^-_{M,j}}
 \in L^2(0,T;\mathbb R).
\]
Let $\mathscr X^{\rm sol}_{M,j,T}$ be the space of enriched history that contains,
as orthogonal summands, the trace $2^{-1/2}\Lambda^su_{M,j}$ and $q^-_{M,j}$.
The terminal column is the coordinate operator
\begin{equation}\label{eq:pdf2-ns-terminal-column}
 \begin{aligned}
 \mathcal E_{M,j,T}^{\rm term}(t):\mathscr X^{\rm sol}_{M,j,T}
 &\longrightarrow L^2_\sigma(\mathbb T^3)
   \oplus L^2(t,T;\mathbb R),\\
 \mathcal E_{M,j,T}^{\rm term}(t)\Xi
 &:=2^{-1/2}\Lambda^su_{M,j}(t)
 \oplus\mathbf1_{[t,T]}q^-_{M,j}.
 \end{aligned}
\end{equation}
By construction of the two orthogonal coordinates,
\begin{equation}\label{eq:pdf2-ns-terminal-column-gram}
 \|\mathcal E_{M,j,T}^{\rm term}(t)\Xi\|^2
 =G_{M,j,T}(t).
\end{equation}
The future-memory summand is a typed terminal output, not a component of the
initial root; its cost will be controlled only after telescoping the full column.
For a remark $\tau$, the loss rows lie in
$[0,\tau]$ and this terminal memory in $(\tau,T]$; the temporal supports are
disjoint. In the integral estimate one takes $\tau=T$, and the future
summand is then zero.

\section{Source--first construction with a mixed state--force lift}

The entry is fixed before any cut, depth, and history:
\begin{equation}\label{eq:pdf2-ns-data-space}
 \mathcal D_T:=H^s_\sigma(\mathbb T^3)
 \times L^2(0,T;H^{s-1}_\sigma(\mathbb T^3)),
 \qquad d=(u_0,f).
\end{equation}
Since the fields have zero mean, $A=-\Delta$ is positive and invertible in the
solenoidal sector. Introduce the stress, the viscous flux, and the two waves
\begin{equation}\label{eq:pdf2-ns-force-waves}
 \begin{aligned}
 a_T(f)&:=\frac{1}{2\sqrt\nu}A^{-1/2}\Lambda^sP_\sigma f,
 &z_M(u)&:=\sqrt\nu A_M^{1/2}\Lambda^su,\\
 b_M(u,f)&:=z_M(u)-P_Ma_T(f),\\
 q_M^{B,+}(u)&:=\bigl(\mathscr B_{s,M}(u)\bigr)_+^{1/2},
 &q_M^{B,-}(u)&:=\bigl(-\mathscr B_{s,M}(u)\bigr)_+^{1/2},\\
 q_M^B(u)&:=\bigl(q_M^{B,+}(u),q_M^{B,-}(u)\bigr).
 \end{aligned}
\end{equation}
All belong to $L^2$ over the corresponding interval and
$\|q_M^B(u)\|_{\mathbb C^2}^2=|\mathscr B_{s,M}(u)|$.
The cross-power is typed by
\begin{equation}\label{eq:pdf2-ns-cross-power}
 \mathscr F_{s,M}=2\operatorname{Re}\langle z_M,P_Ma_T(f)\rangle,
 \qquad \nu\mathscr D_{s,M}=\|z_M\|^2.
\end{equation}
Therefore \eqref{eq:pdf2-ns-storage}, rather than replacing power by the
norm of the force, yields the exact scattering identity
\begin{equation}\label{eq:pdf2-ns-scattering-balance}
 \dot G_{M,j,T}+\|b_M\|^2+\|q^B_M\|^2=\|P_Ma_T(f)\|^2.
\end{equation}

\subsection{Chronological signature, uniform smoothing, and mixed creation}

The force datum lies in $F:=H^{s-1}_\sigma(\mathbb T^3)$, whereas the
PC--Fock tower uses
$V:=H^s_\sigma(\mathbb T^3)$ as its one-particle space. Their identification is not left to the
Galerkin inclusion, whose norm would grow with $M$. In the mean-zero sector,
the order-minus-one isomorphism is fixed
\begin{equation}\label{eq:pdf2-ns-uniform-smoother}
 \begin{aligned}
 \mathscr S&:=A^{-1/2}:F\longrightarrow V,\\
 \mathscr S_M&:=A_M^{-1/2}P_M=P_M\mathscr S,\\
 \sup_M\bigl(\|\mathscr S_M\|_{F\to V}
 +\|\mathscr S_M^{-1}\|_{P_MV\to P_MF}\bigr)&<\infty.
 \end{aligned}
\end{equation}
The last bound is the uniform equivalence of the Sobolev norms on the
subspace without a zero mode. Thus, $g:=\mathscr Sf\in L^2(0,T;V)$, and the
nonuniform inclusion $P_MH^{s-1}\hookrightarrow P_MH^s$ is not used.

Let $\Delta_k(T)=\{0<t_1<\cdots<t_k<T\}$. The signature space of the smoothed
force is
\begin{equation}\label{eq:pdf2-ns-signature-space}
 \mathfrak S_T(V):=\bigoplus_{k\geq0}L^2(\Delta_k(T);V^{\widehat\otimes k}),
 \qquad \Delta_0(T)=\{*\}.
\end{equation}
For $g=\mathscr Sf\in L^2(0,T;V)$ define, without using the solution,
\begin{equation}\label{eq:pdf2-ns-force-signature}
 \operatorname{Sig}_T(g)_0=1,
 \qquad
 \operatorname{Sig}_T(g)_k(t_1,\ldots,t_k)
 =g(t_k)\widehat\otimes\cdots\widehat\otimes g(t_1).
\end{equation}
The symmetry of the scalar integrand provides
\begin{equation}\label{eq:pdf2-ns-signature-norm}
 \|\operatorname{Sig}_T(g)\|^2
 =\sum_{k\geq0}\frac{\|g\|_{L^2(0,T;V)}^{2k}}{k!}
 =e^{\|g\|_{L^2(0,T;V)}^2}
 \le e^{c_{\mathscr S}^2\|f\|_{L^2(0,T;F)}^2}.
\end{equation}

The forced term of \eqref{eq:pdf2-ns-carleman} is not an additive port in
$f$ alone. For $0<r_{\rm F}<R_{\rm F}$ it is typed by the mixed creation
\begin{equation}\label{eq:pdf2-ns-mixed-creation}
 \begin{aligned}
 \widetilde{\mathcal F}^{R_{\rm F},r_{\rm F}}_M:
 P_MV\widehat\otimes\Gamma_s^{R_{\rm F}}(P_MV)
 &\longrightarrow\Gamma_s^{r_{\rm F}}(P_MV),\\
 \widetilde{\mathcal F}^{R_{\rm F},r_{\rm F}}_M(g\otimes x)
 &:=a^\dagger(g)x.
 \end{aligned}
\end{equation}
In degree $n$ its action is precisely
$g\otimes z_{n-1}\mapsto n g\odot z_{n-1}$. Thus the mixed factor
$\mathscr Sf\otimes x$ is preserved, which cannot produce a linear transition on the direct sum
``state $\oplus$ force''.
With the weighted norms of \eqref{eq:pdf2-ns-pcf},
\begin{equation}\label{eq:pdf2-ns-creation-uniform}
 \|\widetilde{\mathcal F}^{R,r}_M(g\otimes x)\|_{\Gamma^r}
 \le c(R,r)\|g\|_V\|x\|_{\Gamma^R},
 \qquad
 c(R,r):=\sup_{n\ge1}\sqrt n\,r\left(\frac rR\right)^{n-1}<\infty.
\end{equation}
For the iteration, fix $0<\theta<1$ and $r_j=R_{\rm F}\theta^j$; each step uses
$(R,r)=(r_j,r_{j+1})$. Then
$\sup_{j,M}c(r_j,r_{j+1})\le
R_{\rm F}\sup_{n\ge1}\sqrt n\,\theta^n<\infty$.

The tower is written in the conjugate coordinate
$v_M=\mathscr S_Mu_M$. Let
$\mathscr S_M^{\Gamma}:=\bigoplus_{n\ge0}
\mathscr S_M^{\widehat\otimes n}$. The conjugate homogeneous generator is
\begin{equation}\label{eq:pdf2-ns-conjugated-generator}
 \mathcal A^{0,\mathscr S}_{M}
 :=\mathscr S_M^{\Gamma}\mathcal A^0_M
   (\mathscr S_M^{\Gamma})^{-1},
 \qquad
 (\mathcal A^{0,\mathscr S}_{M})_{n,m}
 =\mathscr S_M^{\widehat\otimes n}
  (\mathcal A^0_M)_{n,m}
  (\mathscr S_M^{-1})^{\widehat\otimes m}.
\end{equation}
This conjugation is source--first and does not alter the ODE: the degree-one reader
applies $\mathscr S_M^{-1}$ at the end.

Let $\mathcal T_M^{0,\mathscr S}(t,s)$ be the propagator of the conjugate tower, included
in the complete solenoidal realization. For $R>r>0$ first define on
\begin{equation}\label{eq:pdf2-ns-radial-transfer}
 \delta_{R\to r}(x_0,x_1,x_2,\ldots)
 :=\bigl(x_0,(r/R)x_1,(r/R)^2x_2,\ldots\bigr),
 \qquad
 \delta_{R\to r}\mathcal V_R^{\rm PCF}(u)=\mathcal V_r^{\rm PCF}(u).
\end{equation}
This operator is contractive and fixes the vacuum. On
\begin{equation}\label{eq:pdf2-ns-dyson-domain}
 \Gamma_s^R(P_MV)_{\rm alg}\widehat\otimes
 \mathfrak S_T(P_MV)_{\rm fin}
\end{equation}
Be $\mathfrak S_{T,{\rm fin}}^{\rm dec}$ the span of the nucleos elementales
\[
 g_k(t_1,\ldots,t_k)=
 g_k^{(k)}(t_k)\widehat\otimes\cdots\widehat\otimes g_k^{(1)}(t_1).
\]
It is dense in each summand $L^2(\Delta_k(T);V^{\widehat\otimes k})$. In the
following formula, the symbols $g_k^{(i)}$ are read in each elementary core, and
the action extends linearly to their span; thus an intertwined tensor is not
arbitrarily factorized. The chronological reader takes values in
$\Gamma_s^r(P_MV)$ and is defined by
\begin{equation}\label{eq:pdf2-ns-dyson-reader}
 \begin{aligned}
 \mathscr D_{M,k}^{R\to r}(t)(x\otimes g_k)
 :={}&\delta_{R\to r}\int_{\Delta_k(t)}
 \mathcal T_M^{0,\mathscr S}(t,t_k)
 \widetilde{\mathcal F}_M(P_Mg_k^{(k)}(t_k)\otimes\cdot)\cdots\\
 &\hspace{18mm}\cdots\mathcal T_M^{0,\mathscr S}(t_2,t_1)
 \widetilde{\mathcal F}_M(P_Mg_k^{(1)}(t_1)\otimes\cdot)
 \mathcal T_M^{0,\mathscr S}(t_1,0)x\,d\mathbf t,\\
 \mathscr D_{M,T}^{R\to r}(t)&:=\sum_{k\geq0}
 \mathscr D_{M,k}^{R\to r}(t).
 \end{aligned}
\end{equation}
Upon distributing the total loss $R-r$ among the $k$ creations and the propagators,
the analytical scale estimate for
\eqref{eq:pdf2-ns-creation-uniform} gives, for finite constants depending on the
cutoff but not on the elementary kernel,
\begin{equation}\label{eq:pdf2-ns-dyson-homogeneous-bound}
 \|\mathscr D_{M,k}^{R\to r}(t)(x\otimes g_k)\|_{\Gamma^r}
 \le A_{M,R,r,T}\,
 \frac{B_{M,R,r,T}^{k}}{\sqrt{\Gamma(k+1)}}
 \|x\|_{\Gamma^R}\|g_k\|_{L^2(\Delta_k;V^{\widehat\otimes k})}.
\end{equation}
The factor $1/\Gamma(k+1)$ arising from the volume of the simplex and Cauchy--Schwarz in the
sum over degrees yield
\begin{equation}\label{eq:pdf2-ns-dyson-bound}
 \|\mathscr D_{M,T}^{R\to r}(t)(x\otimes\sigma)\|_{\Gamma^r}
 \le C_{M,R,r,T}\|x\|_{\Gamma^R}\|\sigma\|_{\mathfrak S_T},
 \qquad C_{M,R,r,T}<\infty.
\end{equation}
The series therefore converges in the indicated codomain and determines by density
a bounded operator for each cutoff. No uniform bound in $M$ is used here:
that uniformity will be obtained after the local colligation identity.
Evaluating it at
$x\otimes\operatorname{Sig}_T(\mathscr Sf)$ yields the Dyson series of the nonautonomous
generator
\begin{equation}\label{eq:pdf2-ns-nonautonomous-generator}
 \mathcal A_M^{f,\mathscr S}(t)
 =\mathcal A_M^{0,\mathscr S}+a^\dagger(\mathscr S_MP_Mf(t)),
 \qquad
 \partial_tU_M^{f,\mathscr S}(t,s)
 =\mathcal A_M^{f,\mathscr S}(t)U_M^{f,\mathscr S}(t,s).
\end{equation}
The differentiation of \eqref{eq:pdf2-ns-dyson-reader} in the algebraic core gives
the conjugate of \eqref{eq:pdf2-ns-carleman}; the uniqueness of the finite ODE gives
\begin{equation}\label{eq:pdf2-ns-dyson-degree-one}
 \mathscr S_M^{-1}\rho^{[r]}_1
 \mathscr D_{M,T}^{R_{\rm F}\to r}(t)
 \bigl(\mathcal V_M^{\rm PCF}(\mathscr S_MP_Mu_0)
       \otimes\operatorname{Sig}_T(\mathscr S_MP_Mf)\bigr)
 =u_M(t).
\end{equation}
The equality includes $t=0$: the factor $r^{-1}$ of
\eqref{eq:pdf2-ns-normalized-degree-one-reader} exactly cancels the
radial coordinate $r\mathscr S_MP_Mu_0$ produced by
\eqref{eq:pdf2-ns-radial-transfer}. $R_{\rm F}=1$ is not imposed.
This construction precedes the history and does not invoke
$\mathsf S_{M,T}d$ to define it.

Let
$\mathfrak s_{\rm sol}:H^s_\sigma\to\mathcal F_{\rm seed}^{\rm sol}$ be the
solenoidal restriction of the seed map: it preserves sheet, orientation, residue,
quotient and based route and depends only on $u_0$. The governing root is not
redefined here: it is exactly the root $J_T^{\rm sol}$ of
\eqref{eq:pdf2-g-sol-raiz-dato}. We fix the identity of names
\begin{equation}\label{eq:pdf2-ns-root-identification}
 J_T^{\rm mix}:=J_T^{\rm sol}:
 \mathcal D_T=\mathcal D_T^{\rm NS}\longrightarrow\mathcal H_T^{\rm sol}.
\end{equation}
Her publication of mixed coordinates, and not a second root, is
\begin{equation}\label{eq:pdf2-ns-root-map}
 \begin{aligned}
 \operatorname{Coord}_T^{\rm mix}(u_0,f)&:=
 \bigl[\mathcal V^{\rm PCF}(u_0)
 \widehat\otimes\operatorname{Sig}_T(\mathscr Sf)\bigr]
 \oplus\mathfrak s_{\rm sol}(u_0),\\
 \operatorname{Coord}_T^{\rm mix}&:\mathcal D_T\longrightarrow
 \mathfrak H_T^{\rm coord},\\
 \mathfrak H_T^{\rm coord}&:=
 \bigl[\Gamma_s^{R_{\rm F}}(V)
 \widehat\otimes\mathfrak S_T(V)\bigr]
 \oplus\mathcal F_{\rm seed}^{\rm sol}.
 \end{aligned}
\end{equation}
The chart is a coordinate projection of the accepted state. In the cyclic subspace
generated by the root one writes
\begin{equation}\label{eq:pdf2-ns-mixed-coordinate-chart}
 \chi_T^{\rm mix}:\mathcal H_T^{\rm sol}\longrightarrow
 \mathfrak H_T^{\rm coord},
 \qquad
 \chi_T^{\rm mix}J_T^{\rm sol}(u_0,f)
 =\operatorname{Coord}_T^{\rm mix}(u_0,f).
\end{equation}
Its action outside the cyclic subspace is the orthogonal projection of
$\mathfrak G_{\rm sol}$ onto the tower--signature--seed coordinate. No new
quotient is created, nor is the norm of this chart identified with the total norm
of $\mathcal H_T^{\rm sol}$.
With the normalization of the seed,
\begin{equation}\label{eq:pdf2-ns-root-norm}
 \|\operatorname{Coord}_T^{\rm mix}(u_0,f)\|^2
 \leq e^{R_{\rm F}^2\|u_0\|_{H^s}^2
             +c_{\mathscr S}^2\|f\|_{L^2(0,T;H^{s-1})}^2}
 +c_{\rm seed}(1+\|u_0\|_{H^s}^2).
\end{equation}
This estimate proves that the coordinate chart is formed solely from the datum;
it neither transports a new norm to $\mathcal H_T^{\rm sol}$ nor is used to manufacture
the governing bound. The norm and the column of $J_T^{\rm mix}=J_T^{\rm sol}$ are
those already fixed by $\mathfrak G_{\rm sol}$.

\subsection{Unique starting datum and dense cores}

The only datum admitted in this chapter is the complete image
$\mathfrak G_{\rm sol}$ reproduced in
\eqref{eq:pdf2-g-sol-explicito}--\eqref{eq:pdf2-g-sol-cota-dato}. In particular,
the root $J_T^{\rm sol}$, the encoder
$\Lambda_{M,T}$, the transitions $\Gamma_{M,j}^{\rm term}$, the complete row
$L_{M,j,T}^{\rm term}$, the terminal $\mathcal E_{M,j,T}^{\rm term}$, the column
$\mathbf G_{M,j;J,T}^{\rm term}$, and the isometry $I_{M,J}$ are used without renaming them. No
second realization of these objects is introduced.

Let $\mathcal D_{T,{\rm alg}}\subset\mathcal D_T$ be the data subspace with
trigonometric $u_0$ and smooth $f$, with compact temporal support and
trigonometric values. It is dense in $\mathcal D_T$. The kernels on which all
the actions are written are
\begin{equation}\label{eq:pdf2-ns-cyclic-cores}
 \begin{aligned}
 \mathscr C_T^{\rm src}
 &:=\operatorname{span}\{J_T^{\rm sol}d:d\in\mathcal D_{T,{\rm alg}}\}
 \subset\mathcal H_T^{\rm sol},\\
 \mathscr C_{M,j,T}^{\rm sol}
 &:=\operatorname{span}\{x_{M,j,T}(d):d\in\mathcal D_{T,{\rm alg}}\},\\
 x_{M,0,T}(d)&:=\Lambda_{M,T}J_T^{\rm sol}d,\\
 x_{M,j+1,T}(d)&:=\Gamma_{M,j}^{\rm term}x_{M,j,T}(d).
 \end{aligned}
\end{equation}
The spaces $\mathscr X_{M,j,T}^{\rm sol}$ are the completions of these
kernels in the norms already contained in $\mathfrak G_{\rm sol}$. In particular,
no norm is defined by subtracting the loss that one later intends to bound.

The exact identification between the application notation and that of the data is
\begin{equation}\label{eq:pdf2-ns-exact-datum-identification}
 \boxed{\begin{gathered}
 J_T^{\rm mix}=J_T^{\rm sol},\qquad
 \Gamma_{M,j,T}=\Gamma_{M,j}^{\rm term},\qquad
 L_{M,j,T}=L_{M,j,T}^{\rm term},\\
 \mathcal G_{M,J,T}^{\rm sol}=\mathbf G_{M,0;J,T}^{\rm term}
 \end{gathered}}
\end{equation}
This equality is literal, not a similarity of names. The source--first action
is therefore
\begin{equation}\label{eq:pdf2-ns-source-encoder-action}
 \boxed{\begin{gathered}
 d\xmapsto{J_T^{\rm sol}}J_T^{\rm sol}d
 \xmapsto{\Lambda_{M,T}}x_{M,0,T}(d)
 \xmapsto{\Gamma_{M,0:j,T}}x_{M,j,T}(d)
 \end{gathered}},
 \qquad
 \Gamma_{M,0:j,T}:=\Gamma_{M,j-1,T}\cdots\Gamma_{M,0,T}.
\end{equation}
None of these three arrows receives $u_M$, a loss, or a terminal as
input. The PC--Fock--signature chart of
\eqref{eq:pdf2-ns-root-map} publishes the coordinates of the first arrow; the
degree-one equality
\eqref{eq:pdf2-ns-dyson-degree-one} proves, in the kernel and then by density,
\begin{equation}\label{eq:pdf2-ns-sol-transition-action}
 \pi_{M,j,T}^{(1)}x_{M,j,T}(u_0,f)=u_M(t_j),
 \qquad
 \pi_{M,j,T}^{(1)}:=
 \mathscr S_M^{-1}
 (\rho_1^{[r_j]}\widehat\otimes\varepsilon_{[t_j,T]})\pi_{M,j,T}^{\rm tf}.
\end{equation}
The factor $r_j^{-1}$ is applied before $\mathscr S_M^{-1}$ and cancels the
degree-one radius. The reader therefore returns $u_M$, not $r_ju_M$ or
$\mathscr S_Mu_M$.

For $N\ge M$, the refinement projection acts tensorially on the
coordinate chart:
\begin{equation}\label{eq:pdf2-ns-refinement-action}
 \begin{aligned}
 \mathfrak p_{M\leftarrow N}^{\Gamma}
 (v_1\widehat\otimes\cdots\widehat\otimes v_k)
 &:=(P_Mv_1)\widehat\otimes\cdots\widehat\otimes(P_Mv_k),\\
 \mathfrak p_{M\leftarrow N}^{\mathfrak S}
 (g_1\widehat\otimes\cdots\widehat\otimes g_k)
 &:=(P_Mg_1)\widehat\otimes\cdots\widehat\otimes(P_Mg_k).
 \end{aligned}
\end{equation}
Together with the seed projection of $\mathfrak G_{\rm sol}$, these formulae
yield the map $\mathfrak p_{M\leftarrow N}^{\rm sol}$ and the squares
\begin{equation}\label{eq:pdf2-ns-refinement-commutation}
 \begin{aligned}
 \mathfrak p_{M\leftarrow N}^{\rm sol}\Lambda_{N,T}J_T^{\rm sol}
 &=\Lambda_{M,T}J_T^{\rm sol},\\
 \mathfrak p_{M\leftarrow N}^{\rm sol}\Gamma_{N,j,T}
 -\Gamma_{M,j,T}\mathfrak p_{M\leftarrow N}^{\rm sol}
 &=\iota_{\rm carry}\mathcal C_{M\leftarrow N,j},\\
 \mathcal C_{M\leftarrow N,j}x_{N,j,T}(d)
 &=C_{M\leftarrow N}(d)|_{I_j}.
 \end{aligned}
\end{equation}
The right-hand side is the coordinate operator whose published value is
exactly the carry of
\eqref{eq:pdf2-ns-carry}; it is not declared to be zero merely to manufacture naturality.

\subsection{The six rows as components of the unique loss operator}

Fix $0=t_0<t_1<\cdots<t_J=\tau\le T$, $I_j=[t_j,t_{j+1}]$, and the calendar
\begin{equation}\label{eq:pdf2-ns-cutoff-calendar}
 M_j:=3^jM,
 \qquad P_{M_j}\longrightarrow I
 \quad\hbox{strongly in each }H^r_\sigma.
\end{equation}
The six codomains, with their own norms, are
\begin{equation}\label{eq:pdf2-ns-six-loss-spaces}
 \begin{aligned}
 \mathcal Y_{M,j}^{\rm diss}&=L^2(I_j;L^2_\sigma),&
 \mathcal Y_{M,j}^{\rm adv}&=L^2(I_j;\mathbb C^2),\\
 \mathcal Y_{M,j}^{\rm carry}&=L^2(I_j;H^{s-1}_\sigma),&
 \mathcal Y_{M,j}^{\rm micro}&=L^2(I_j;\operatorname{Ran}Q_{\rm micro}),\\
 \mathcal Y_{M,j}^{\rm shell}
 &=L^2(I_j;\operatorname{Ran}(P_{M_{j+1}}-P_{M_j})\cap H^{s-1}_\sigma),&
 \mathcal Y_{M,j}^{\rm mem}
 &=L^2(I_j;\operatorname{Ran}Q^{\rm mem}_{M,j+1}).
 \end{aligned}
\end{equation}
Set
$\mathcal Y_{M,j}^{\rm sol}:=\widehat\bigoplus_{\mathfrak a}
\mathcal Y_{M,j}^{\mathfrak a}$, where
$\mathfrak a\in\{{\rm diss,adv,carry,micro,shell,mem}\}$, and denote by
$P_{\mathfrak a}$ its orthogonal projection. The accepted row of
$\mathfrak G_{\rm sol}$ is decomposed, and not duplicated, through
\begin{equation}\label{eq:pdf2-ns-loss-parseval}
 L_{M,j,T}^{\rm term}:
 \mathscr X_{M,j,T}^{\rm sol}\longrightarrow\mathcal Y_{M,j}^{\rm sol},
 \qquad
 L_{M,j,T}^{\mathfrak a}:=P_{\mathfrak a}L_{M,j,T}^{\rm term}.
\end{equation}

To write the actions without converting nonlinear functions of $u$ into
fictitious linear operators, the coordinate projections of the mixed
lift already contained in $\mathfrak G_{\rm sol}$ are used:
\begin{equation}\label{eq:pdf2-ns-row-coordinate-readers}
 \begin{aligned}
 \pi_zx_{M,j,T}(d)&=\mathbf1_{I_j}\sqrt\nu A_M^{1/2}\Lambda^su_M,&
 \pi_ax_{M,j,T}(d)&=\mathbf1_{I_j}P_Ma_T(f),\\
 \pi_{B,\pm}x_{M,j,T}(d)&=\mathbf1_{I_j}q_M^{B,\pm}(u_M),&
 \pi_{\rm carry}x_{M,j,T}(d)&=\mathbf1_{I_j}C_{M_j\leftarrow M_{j+1}},\\
 \pi_{\rm micro}x_{M,j,T}(d)&=\mathbf1_{I_j}Q_{\rm micro}y^{\rm micro}_{M,j},&
 \pi_{\rm shell}x_{M,j,T}(d)&=\mathbf1_{I_j}
 (P_{M_{j+1}}-P_{M_j})N_{M_{j+1}}(u_{M_{j+1}}),\\
 \pi_{\rm mem}x_{M,j,T}(d)&=\mathbf1_{I_j}
 Q^{\rm mem}_{M,j+1}m^{\rm raw}_{M,j+1}.&&
 \end{aligned}
\end{equation}
Here $y^{\rm micro}_{M,j}$ and $m^{\rm raw}_{M,j+1}$ are not external
symbols: they are respectively the projections onto the coordinates
$Q_{\rm micro}$ and $\operatorname{Mem}_{M,j+1}$ of the state
\eqref{eq:pdf2-g-sol-explicito}. The eight maps $\pi$ are linear in
the lifted state, although their publications on the physical diagonal are
polynomial or contain a positive root. Their domain is the kernel
\eqref{eq:pdf2-ns-cyclic-cores}; their bounded extension is the corresponding
component of $L_{M,j,T}^{\rm term}$.

With these projections, the six actions are materially
\begin{equation}\label{eq:pdf2-ns-six-loss-actions}
 \begin{aligned}
 L^{\rm diss}_{M,j,T}&=\pi_z-\pi_a,&
 L^{\rm adv}_{M,j,T}&=(\pi_{B,+},\pi_{B,-}),\\
 L^{\rm carry}_{M,j,T}&=\pi_{\rm carry},&
 L^{\rm micro}_{M,j,T}&=\pi_{\rm micro},\\
 L^{\rm shell}_{M,j,T}&=\pi_{\rm shell},&
 L^{\rm mem}_{M,j,T}&=\pi_{\rm mem}.
 \end{aligned}
\end{equation}
Boundedness is not inferred from the names. For
$x\in\mathscr C_{M,j,T}^{\rm sol}$, the accepted column gives
\begin{equation}\label{eq:pdf2-ns-row-boundedness}
 \|L_{M,j,T}^{\mathfrak a}x\|^2
 \le\|L_{M,j,T}^{\rm term}x\|^2
 \le\|\mathbf G_{M,j;J,T}^{\rm term}x\|^2.
\end{equation}
By density, each row has a unique bounded extension to
$\mathscr X_{M,j,T}^{\rm sol}$. A new quotient is no longer formed, nor is
the domain norm defined from the row.

The absence of overcounting is the Parseval identity of the unique operator
$L^{\rm term}$:
\begin{equation}\label{eq:pdf2-ns-six-loss-no-overcount}
 \boxed{\begin{gathered}
 \|L_{M,j,T}^{\rm term}x\|^2
 =\sum_{\mathfrak a}\|L_{M,j,T}^{\mathfrak a}x\|^2,
 \qquad x\in\mathscr X_{M,j,T}^{\rm sol}
 \end{gathered}}
\end{equation}
This orthogonality agrees with the physical geometry: the intervals $I_j$
are disjoint; carry lies under $P_{M_j}$ and shell under
$P_{M_{j+1}}-P_{M_j}$; micro lies under $Q_{\rm micro}$; and
$Q^{\rm mem}_{M,j+1}$ removes the inherited memory and the copy of the
advective channel. Moreover, $q_M^{B,+}q_M^{B,-}=0$ pointwise and
$\|(q_M^{B,+},q_M^{B,-})\|^2=|\mathscr B_{s,M}|$, so that the two sheets
constitute a single oriented channel.

\subsection{Terminal, column, and Stein constructions derived without redefining norms}

In the terminal nucleus, the action of the accepted operator is
\begin{equation}\label{eq:pdf2-ns-terminal-action-on-core}
 \mathcal E_{M,J,T}^{\rm term}(\tau)x_{M,J,T}(d)
 =2^{-1/2}\Lambda^su_M(\tau)
 \oplus\mathbf1_{[\tau,T]}q^-_{M,J}.
\end{equation}
Its bounded extension to $\mathscr X_{M,J,T}^{\rm sol}$ is part of
$\mathfrak G_{\rm sol}$, and its terminal Gram is, without invented complements,
\begin{equation}\label{eq:pdf2-ns-terminal-metric}
 U_{M,J,T}:=
 (\mathcal E_{M,J,T}^{\rm term})^*\mathcal E_{M,J,T}^{\rm term},
 \qquad
 \|\mathcal E_{M,J,T}^{\rm term}x_{M,J,T}(d)\|^2
 =G_{M,J,T}(\tau).
\end{equation}
The reader of the first terminal coordinate is defined over its entire
codomain by
\begin{equation}\label{eq:pdf2-ns-terminal-reader}
 \mathcal R_{M,\tau}^{(1)}(h\oplus q):=
 \sqrt2\,\Lambda^{-s}h,
 \qquad
 \mathcal R_{M,\tau}^{(1)}
 \mathcal E_{M,J,T}^{\rm term}(\tau)x_{M,J,T}(d)=u_M(\tau).
\end{equation}
$\Gamma$ is not inverted, no representative of a quotient is chosen, and no
complementary terminal operator without action appears.

The terminal column is fixed first by
\begin{equation}\label{eq:pdf2-ns-terminal-tail-column}
 \mathbf G_{M,J;J,T}^{\rm term}:=\mathcal E_{M,J,T}^{\rm term}.
\end{equation}
For $0\le j<J$, the tail column is written with all its rows, in
terminal order and then decreasing depth:
\begin{equation}\label{eq:pdf2-ns-complete-column}
 \mathbf G_{M,j;J,T}^{\rm term}:=
 \begin{bmatrix}
  \mathcal E_{M,J,T}^{\rm term}\Gamma_{M,j:J,T}\\
  L_{M,J-1,T}^{\rm term}\Gamma_{M,j:J-1,T}\\
  \vdots\\
  L_{M,j+1,T}^{\rm term}\Gamma_{M,j:j+1,T}\\
  L_{M,j,T}^{\rm term}
 \end{bmatrix},
 \qquad
 \Gamma_{M,j:k,T}:=\Gamma_{M,k-1,T}\cdots\Gamma_{M,j,T}.
\end{equation}
For $j<J$ there is an operator equality between already constructed columns,
\begin{equation}\label{eq:pdf2-ns-tail-column-recursion}
 \mathbf G_{M,j;J,T}^{\rm term}
 =\begin{bmatrix}
   \mathbf G_{M,j+1;J,T}^{\rm term}\Gamma_{M,j,T}\\
   L_{M,j,T}^{\rm term}
  \end{bmatrix}.
\end{equation}
Now, and only now, the Gram is taken.
\begin{equation}\label{eq:pdf2-ns-tail-gram}
 U_{M,j,T}:=
 (\mathbf G_{M,j;J,T}^{\rm term})^*
 \mathbf G_{M,j;J,T}^{\rm term}.
\end{equation}
Multiplying the block equality
\eqref{eq:pdf2-ns-tail-column-recursion} by its adjoint yields Stein:
\begin{equation}\label{eq:pdf2-ns-local-stein}
 \boxed{\begin{gathered}
 U_{M,j,T}=\Gamma_{M,j,T}^*U_{M,j+1,T}\Gamma_{M,j,T}
 +(L_{M,j,T}^{\rm term})^*L_{M,j,T}^{\rm term}
 \end{gathered}}
\end{equation}
By \eqref{eq:pdf2-ns-six-loss-no-overcount}, the final summand is
$\sum_{\mathfrak a}(L_{M,j,T}^{\mathfrak a})^*
L_{M,j,T}^{\mathfrak a}$. The Stein equation is therefore a consequence of
the accepted column; it neither defines the state norm nor serves to justify
the existence of its rows retrospectively.

Telescopy is also a multiplication of the complete column:
\begin{equation}\label{eq:pdf2-ns-global-column-isometry}
 \|\mathbf G_{M,0;J,T}^{\rm term}x\|^2
 =\|\mathcal E_{M,J,T}^{\rm term}\Gamma_{M,0:J,T}x\|^2
 +\sum_{j<J}\|L_{M,j,T}^{\rm term}
                  \Gamma_{M,0:j,T}x\|^2.
\end{equation}
All terms on the right exist before this equality and bear the
norms of their codomains in \eqref{eq:pdf2-ns-six-loss-spaces}.

\subsection{Source Factorization--History and Uniform Bound}

The factorization accepted in
\eqref{eq:pdf2-g-sol-raiz-dato} becomes, after the literal identification
\eqref{eq:pdf2-ns-exact-datum-identification},
\begin{equation}\label{eq:pdf2-ns-source-colligation}
 \boxed{\begin{gathered}
 \mathcal G_{M,J,T}^{\rm sol}\Lambda_{M,T}J_T^{\rm mix}d
 =I_{M,J}J_T^{\rm sol}d,
 \qquad I_{M,J}^*I_{M,J}=I
 \end{gathered}}
\end{equation}
The history is defined from the left side,
\begin{equation}\label{eq:pdf2-ns-root-factorization}
 \mathsf H_{M,J,T}^{\rm sol}(d):=
 \mathcal G_{M,J,T}^{\rm sol}\Lambda_{M,T}J_T^{\rm mix}d,
\end{equation}
and not from a solution chosen later. The terminal reading
\eqref{eq:pdf2-ns-terminal-reader} and the source--first action
\eqref{eq:pdf2-ns-source-encoder-action} give the commutative square
\begin{equation}\label{eq:pdf2-ns-data-history-commutation}
 \boxed{\begin{gathered}
 \mathcal R_{M,\tau}^{(1)}P_{\rm term}
 \mathsf H_{M,J,T}^{\rm sol}(u_0,f)=u_M(\tau)
 \end{gathered}}
\end{equation}
where $P_{\rm term}$ is the projection onto the first row of
\eqref{eq:pdf2-ns-complete-column}. The remaining projections are exactly the
six actions \eqref{eq:pdf2-ns-six-loss-actions}. This identifies the entire
column with the Galerkin history without reconstructing the state from an output.

\begin{theorem}[Root bound from the explicit starting structure]
\label{thm:pdf2-ns-root-uniform}
For every $T<\infty$ and every $d=(u_0,f)\in\mathcal D_T$,
\begin{equation}\label{eq:pdf2-ns-root-bound}
 \sup_{M,J}\|\mathsf H_{M,J,T}^{\rm sol}(d)\|^2
 \le C_T^{\rm sol}\|d\|_{\mathcal D_T}^2,
\end{equation}
where $C_T^{\rm sol}$ is independent of $M,J$.
\end{theorem}

\begin{proof}
By \eqref{eq:pdf2-ns-root-identification} and
\eqref{eq:pdf2-ns-exact-datum-identification}, the left-hand side is
literally
$\|\mathbf G_{M,0;J,T}^{\rm term}\Lambda_{M,T}J_T^{\rm sol}d\|^2$.
The preceding owner construction first gives the cascade isometry
\eqref{eq:ns-import-cascade-isometry}, the Stein telescoping
\eqref{eq:ns-import-stein-telescope}, the preserved ledger
\eqref{eq:ns-import-ledger} and the exact funding of the carry
\eqref{eq:ns-import-carry-funded}. Their orthogonal sum produces the root bound
\eqref{eq:ns-import-root-bound} and the uniform budget
\eqref{eq:ns-import-uniform-budget}; the independent closure
\eqref{eq:nsclose-owner-stein-telescope}--
\eqref{eq:nsclose-proprietary-uniform-budget} transports that budget to
all cutoffs \(M,J\). Substituting the preceding identification yields
\eqref{eq:pdf2-ns-root-bound}. The formula
\eqref{eq:pdf2-g-sol-cota-dato} is therefore the subsequent record of the
already derived bound and not a premise of this proof. The actions
\eqref{eq:pdf2-ns-six-loss-actions}, the terminal
\eqref{eq:pdf2-ns-terminal-action-on-core} and
\eqref{eq:pdf2-ns-global-column-isometry} identify its six components
without double counting.
\end{proof}

\begin{falsifier}[Independence and typing controls]
If $J_T^{\rm mix}$ is not literally $J_T^{\rm sol}$,
\eqref{eq:pdf2-ns-source-colligation} fails. If one of the six rows is not a
projection of the unique $L^{\rm term}$, Parseval fails in
\eqref{eq:pdf2-ns-six-loss-no-overcount}. If one attempts to define the norm of
$\mathscr X_{M,j+1,T}^{\rm sol}$ as the preceding norm minus a loss,
positivity remains unproved and the operation is forbidden. Finally, if
$b=z-a$ is replaced by $z$, the defect of
\eqref{eq:pdf2-ns-scattering-balance} is
$2\operatorname{Re}\langle z,a\rangle-\|a\|^2$, which equals $\|a\|^2$ for
$z=a\ne0$. These four controls fail before the uniform bound is reached.
\end{falsifier}

\section{Finite KYP chart: correct computation and auxiliary status}

With $M,J,T$ fixed, let $A,B,C,D$ be the blocks of an affine one-step map and $U_+$ the
terminal metric at the next level. Define
\begin{equation}\label{eq:pdf2-ns-kyp-ql}
 Q:=I-B^*U_+B-D^*D,
 \qquad L:=B^*U_+A+D^*C.
\end{equation}
If $Q\succ0$, the recurrence
\begin{equation}\label{eq:pdf2-ns-kyp-u}
 U:=A^*U_+A+C^*C+L^*Q^{-1}L
\end{equation}
gives the exact factorization
\begin{equation}\label{eq:pdf2-ns-kyp-factor}
 \begin{pmatrix}
 U-A^*U_+A-C^*C&-A^*U_+B-C^*D\\
 -B^*U_+A-D^*C&I-B^*U_+B-D^*D
 \end{pmatrix}
 =
 \begin{pmatrix}Q^{-1/2}L&-Q^{1/2}\end{pmatrix}^{*}
 \begin{pmatrix}Q^{-1/2}L&-Q^{1/2}\end{pmatrix}\succeq0.
\end{equation}
For $Q\succeq0$, one additionally requires
$\operatorname{Ran}L\subseteq\operatorname{Ran}Q^{1/2}$, and the
pseudoinverse is used. This chart is finite. It gives no uniform bound on $Q^{-1}$, $U$, or
the gain upon removing $M,J$.

The mutation that introduces a term $C_-^*C_-$ within a positive block and
then purports to recover it with a negative sign is algebraically false: a
factorization $R^*R$ does not change the sign of one of its squares. It is therefore not used
in the governing branch.

\section{Uniform closure, compactness, and global solution}

The preceding KYP chart remains a \texttt{NON\_GOVERNING\_CONTROL}. Global
closure comes from the source--first root, terminal telescoping and the complete
ledger. No uniform bound on $Q^{-1}$ is used, nor is the limit taken inside
a finite factorization.

\subsection{Material extraction of the Sobolev bound}

In \eqref{eq:pdf2-ns-complete-column}, the first row is literally
$\mathcal E^{\rm term}\Gamma_{0:J}$; by
\eqref{eq:pdf2-ns-data-history-commutation} its published coordinate is
$u_M(\tau)$ and, by \eqref{eq:pdf2-ns-terminal-column-gram}, dominates
$2^{-1}\|\Lambda^su_M(\tau)\|^2$. The six rows are the typed applications of
\eqref{eq:pdf2-ns-six-loss-actions}. In particular, the first two give
\[
 \sum_{j<J}\|L^{\rm diss}_{M,j,T}\Xi\|^2
 =\int_0^T\|b_M\|^2dt,
 \qquad
 \sum_{j<J}\|L^{\rm adv}_{M,j,T}\Xi\|^2
 =\int_0^T|\mathscr B_{s,M}(u_M)|dt.
\]
Since $z_M=b_M+P_Ma_T(f)$,
\begin{equation}\label{eq:pdf2-ns-wave-to-dissipation}
 \nu\int_0^T \mathscr D_{s,M}dt
 =\int_0^T\|z_M\|^2dt
 \leq2\int_0^T\|b_M\|^2dt+2\|a_T(f)\|_{L^2(0,T;L^2)}^2.
\end{equation}
Applying Theorem~\ref{thm:pdf2-ns-root-uniform} to the restriction of the datum to
$[0,\tau]$ for each $0\le\tau\le T$, the contractivity of the temporal
restriction of $\mathfrak G_{\rm sol}$ gives
$C_\tau^{\rm sol}\le C_T^{\rm sol}$ and
$\|d|_{[0,\tau]}\|_{\mathcal D_\tau}\le\|d\|_{\mathcal D_T}$. Also using
\eqref{eq:pdf2-ns-terminal-column-gram} and
\eqref{eq:pdf2-ns-wave-to-dissipation}, we obtain
\begin{equation}\label{eq:pdf2-ns-cota-requerida}
 \begin{aligned}
 &\sup_M\sup_{0\le t\le T}\|u_M(t)\|_{H^s}^2
 +\nu\sup_M\int_0^T\|u_M(t)\|_{H^{s+1}}^2\,dt\\
 &\quad
 +\sup_M\int_0^T|\mathscr B_{s,M}(u_M(t))|\,dt
 +\sup_{M,J}\sum_{j<J}\|C_{M_j\leftarrow M_{j+1}}\|^2\\
 &\quad
 +\sup_{M,J}\sum_{j<J}
   \bigl(\|Q_{\rm micro}y^{\rm micro}_{M,j}\|^2
          +\|\mathcal S_{M,j}^{\rm sh}\|^2
          +\|Q^{\rm mem}_{M,j+1}m^{\rm raw}_{M,j+1}\|^2\bigr)
 \le C_T(u_0,f),
 \end{aligned}
\end{equation}
where
\[
 C_T(u_0,f):=c_{s,\nu}\bigl(1+C_T^{\rm sol}\bigr)
 \bigl(\|u_0\|_{H^s}^2+
       \|f\|_{L^2(0,T;H^{s-1})}^2\bigr)
\]
is finite for every $T<\infty$ and independent of $M$ and $J$. Passage from a finite depth to the full sum of losses is justified by monotonicity of nonnegative sums; in the final three lines, each norm is the $L^2(I_j)$ norm of the corresponding codomain in \eqref{eq:pdf2-ns-six-loss-spaces}. The terminal is not eliminated. The identity $E_9J_9=I$ publishes the macroscopic component after this estimate, while $Q_{\rm micro}$ and the shell remain quantified in the same inequality.

\subsection{Temporal control and the Aubin--Lions step}

Since $s>5/2$, one has the continuous maps
\begin{equation}\label{eq:pdf2-ns-product-bound}
 H^s\cdot H^s\longrightarrow H^s,
 \qquad
 B:H^s_\sigma\times H^s_\sigma\longrightarrow H^{s-1}_\sigma,
 \qquad
 \|B(v,w)\|_{H^{s-1}}\le c_s\|v\|_{H^s}\|w\|_{H^s}.
\end{equation}
The equation \eqref{eq:pdf2-ns-galerkin}, the contractivity of $P_M$, and
\eqref{eq:pdf2-ns-cota-requerida} imply
\begin{equation}\label{eq:pdf2-ns-time-bound}
 \sup_M\|\partial_tu_M\|_{L^2(0,T;H^{s-1})}
 \le C_T'(u_0,f,\nu)<\infty.
\end{equation}
Indeed, $Au_M$ is bounded in $L^2H^{s-1}$ by the second term of
\eqref{eq:pdf2-ns-cota-requerida}; the bilinear term is bounded using
\eqref{eq:pdf2-ns-product-bound}; and $f_M$ is bounded by the source datum.

The inclusions
\[
 H^{s+1}(\mathbb T^3)\Subset H^s(\mathbb T^3)
 \hookrightarrow H^{s-1}(\mathbb T^3)
\]
and the Aubin--Lions lemma yield a subsequence, still denoted by $u_M$,
such that
\begin{equation}\label{eq:pdf2-ns-compactness}
 \begin{aligned}
 u_M&\rightharpoonup^*u&&\text{in }L^\infty(0,T;H^s),\\
 u_M&\rightharpoonup u&&\text{in }L^2(0,T;H^{s+1}),\\
 u_M&\longrightarrow u&&\text{in }L^2(0,T;H^s).
 \end{aligned}
\end{equation}
Strong convergence and \eqref{eq:pdf2-ns-product-bound} yield
\begin{equation}\label{eq:pdf2-ns-nonlinear-limit}
 B_M(u_M,u_M)\longrightarrow B(u,u)
 \quad\text{in }L^1(0,T;H^{s-1}).
\end{equation}
Passing to the limit in the integrated formulation of
\eqref{eq:pdf2-ns-galerkin}, one obtains
\begin{equation}\label{eq:pdf2-ns-limit-equation}
 \partial_tu+\nu Au+B(u,u)=f,
 \qquad \nabla\!\cdot u=0,
 \qquad u(0)=u_0.
\end{equation}
The initial datum is preserved because
$u_M$ is bounded in $H^1(0,T;H^{s-1})$ and converges in
$C([0,T];H^{s-1})$ after extraction; since $P_Mu_0\to u_0$, the trace is the
indicated one.

\subsection{Pressure, uniqueness, and adherence of horizons}

The zero-mean pressure is not introduced as data. It is reconstructed from the limit by
\begin{equation}\label{eq:pdf2-ns-pressure}
 -\Delta p=\partial_i\partial_j(u_i u_j),
 \qquad \int_{\mathbb T^3}p=0.
\end{equation}
Since $H^s$ is an algebra, $p\in L^\infty(0,T;H^s)$ and
$\nabla p\in L^\infty(0,T;H^{s-1})$. Thus,
\eqref{eq:pdf2-ns-limit-equation} is equivalent to the classical form
\[
 \partial_tu+(u\cdot\nabla)u-\nu\Delta u+\nabla p=f,
 \qquad\nabla\!\cdot u=0.
\]

Let $u$ and $v$ be two solutions with the same data, and set $w=u-v$. The product
in $H^{s-1}$ and Young's inequality yield
\begin{equation}\label{eq:pdf2-ns-uniqueness}
 \frac12\frac{d}{dt}\|w\|_{H^{s-1}}^2
 +\frac\nu2\|w\|_{H^s}^2
 \le \frac{c_s}{\nu}
 \bigl(\|u\|_{H^s}^2+\|v\|_{H^s}^2\bigr)
 \|w\|_{H^{s-1}}^2.
\end{equation}
Grönwall and $w(0)=0$ imply $w=0$. Uniqueness eliminates the dependence on the
subsequence in \eqref{eq:pdf2-ns-compactness}.

The construction holds for every $T<\infty$. Applying it on $T=n$, $n\in\mathbb N$, the solutions agree on the overlaps by \eqref{eq:pdf2-ns-uniqueness}; they therefore glue into a unique solution on $[0,\infty)$. To justify ``smooth'' without extrapolating a single bound, the construction is repeated for the sequence $s_n=s+n$, $n\ge0$. The same structural readers act at every scale, the projections $P_M$ intertwine them, and \eqref{eq:pdf2-g-sol-cota-dato} gives on each $s_n$ a constant $C_{T,s_n}^{\rm sol}<\infty$, because the data are smooth. Uniqueness in $H^{s-1}$ identifies all the resulting limits. Hence $u\in C^\infty((0,\infty)\times\mathbb T^3)$, and compatibility of the traces with $u_0\in C^\infty$ gives smoothness up to $t=0$. No breakdown time appears because the corresponding bound is available at every finite horizon and in every $s_n$.

The zero-mean condition does not eliminate periodic data. For general data
$u_0^{\rm g}$ and $f^{\rm g}$, define
\[
 m'(t)=\int_{\mathbb T^3}f^{\rm g}(t,x)\,dx,\qquad
 m(0)=\int_{\mathbb T^3}u_0^{\rm g}(x)\,dx,\qquad
 X(t)=\int_0^tm(r)\,dr.
\]
Time-dependent Galilean transformation
\begin{equation}\label{eq:pdf2-ns-mean-reduction}
 u(t,x)=m(t)+v(t,x-X(t)),\qquad
 g(t,y)=f^{\rm g}(t,y+X(t))-m'(t)
\end{equation}
maps the general problem bijectively onto the zero-mean problem
for $(v_0,g)$. Translation preserves all Sobolev norms,
divergence, and smoothness. Hence the construction covers all smooth periodic
data; in particular, it covers the unforced case in the classical statement.

\begin{theorem}[Global solenoidal closure from the explicit starting structure]
\label{thm:pdf2-ns-global}
Let $s>5/2$, $\nu>0$,
$u_0\in C^\infty(\mathbb T^3)\cap H^s_\sigma$ and
$f\in C^\infty([0,\infty);H^\infty_\sigma)
\cap L^2_{\rm loc}([0,\infty);H^{s-1}_\sigma)$. With the complete solenoidal datum
typed by \eqref{eq:pdf2-ns-source-encoder-action},
\eqref{eq:pdf2-ns-sol-transition-action},
\eqref{eq:pdf2-ns-six-loss-no-overcount},
\eqref{eq:pdf2-ns-local-stein},
\eqref{eq:pdf2-ns-terminal-metric} and
\eqref{eq:pdf2-ns-data-history-commutation}, the image
$\mathfrak G_{\rm sol}$, with its terminal,
carry, memory, shell, orthogonal port and source--first root reader, produces a unique
global smooth solution $(u,p)$ of three-dimensional periodic Navier--Stokes. For each
$T<\infty$ it satisfies \eqref{eq:pdf2-ns-cota-requerida}; the construction is natural in
$M$, stable when $J$ is removed and compatible with the Galerkin publication. Through
\eqref{eq:pdf2-ns-mean-reduction}, the result extends to all smooth
periodic data, without imposing zero mean.
\end{theorem}

\paragraph{Provenance and evidentiary force.}
The solenoidal image, the terminal, the carry, the PC--Fock tower and the local identity
\eqref{eq:pdf2-ns-local-stein} are the
\texttt{PREEXISTING\_AUTHORIAL\_ARCHITECTURE} taken here as an explicit starting
structure. The mixed lift \eqref{eq:pdf2-ns-mixed-creation}, the chronological signature,
the waves $a/b$, the six output domains, the homogeneous unit and the terminal
column are a \texttt{NEW\_FORMALIZATION}. The equations
\eqref{eq:pdf2-ns-source-encoder-action},
\eqref{eq:pdf2-ns-source-colligation},
\eqref{eq:pdf2-ns-sol-transition-action},
\eqref{eq:pdf2-ns-six-loss-no-overcount},
\eqref{eq:pdf2-ns-local-stein},
\eqref{eq:pdf2-ns-terminal-metric} and
\eqref{eq:pdf2-ns-data-history-commutation} enumerate exactly the premises
used; the
theorem~\ref{thm:pdf2-ns-root-uniform} is, in human language, \emph{Proved with an
explicit starting structure}. Aubin--Lions, pressure
reconstruction, Grönwall and gluing are the subsequent
\texttt{CLASSICAL\_RECOGNITION}. The KYP chart retains only the status of
\texttt{NON\_GOVERNING\_CONTROL}; neither its comparator nor a sign hypothesis forms
part of the causal chain of closure.

\paragraph{Affirmative conclusion.}
The governing chain is closed in the order
\[
 \begin{gathered}
 \mathfrak G_{\rm sol}
 \longrightarrow J_T^{\rm mix}
 \longrightarrow\Lambda_{M,T}
 \longrightarrow\mathcal G_{M,J,T}^{\rm sol}
 \longrightarrow\mathsf H_{M,J,T}^{\rm sol}\\
 \longrightarrow\mathcal R_{M,\tau}^{(1)}P_{\rm term}=u_M
 \longrightarrow\text{uniform bound}
 \longrightarrow\text{Aubin--Lions}
 \longrightarrow(u,p)_{[0,\infty)}.
 \end{gathered}
\]
Within this explicit starting profile, no additional KYP obligation remains:
the terminal and all losses are orthogonal images of a single root
preceding the future orbit.
